% Аутоматски генерисано из ОЗП.docx

\chapter{Откривање законитости у подацима}
\label{ch:ozp}
\textbf{Откривање законитости у подацима (ОЗП)} представља процес у коме се откривају законитости у подацима организација са циљем унапређења и/или аутоматизације процеса одлучивања. Под подацима се подразумевају скупови података, најчешће представљени у табеларној форми (редови табеле представљају случајеве, док колоне представљају атрибуте који описују случајеве), који су добијени из база података, складишта података, али и свих других извора преко којих се прикупљају и складиште подаци (текст, логови, фото, видео и аудио записи).

ОЗП обједињује методе, алгоритме и технике из различитих дисциплина (статистика и вероватноћа, математика, машинско учење, базе података, структуре података). Може да се користи у анализи великих података у организацијама — за разлику од класичне статистичке анализе која је обично користила табеле мањих димензија, ОЗП је намењен за анализу већих количина података и уводи низ нових алгоритама који могу да раде са великим обимом података, разноврснијим типовима података, као и учесталијом фреквенцијом генерисања нових података.

\textbf{Законитости} су правила или модели који важе у специфичном контексту, али не морају нужно да важе у свим ситуацијама (нпр. понашање потрошача у Азији и Европи не мора да подлеже истим законитостима). Са друге стране, \textbf{закони} су оне појаве које имају далеко већи опсег дејства и далеко шири контекст (закон понуде и потражње, закон гравитације). Често се за реч законитост, поготово у инжењерским наукама, користи и појам \textit{патерн}, \textit{архетип}, \textit{узор}.

Откривене законитости су изузетно важне за сваку организацију јер јој омогућавају, између осталог:

\begin{itemize}
  \item боље разумевање пословног процеса,
  \item информисаније доношење одлука,
  \item потенцијално претварање законитости у нову вредност, а самим тим и
  \item могућност аутоматизације процеса одлучивања.
\end{itemize}
Информација (управљачки извештај) оставља простор доносиоцу одлука да донесе одлуку, док откривена законитост може да аутоматизује пословни процес, јер садржи, за разлику од информације, управљачку компоненту.

ОЗП не представља замену за класичну статистичку анализу, већ њену допуну. Знања која су потребна за разумевање метода и резултата статистичке анализе остају и даље неопходна. Истовремено, ма колико примамљиво откривање законитости из података звучи, треба бити опрезан, јер неки кандидати за законитости (идентификована правила или међузависности) могу пре да буду окарактерисани као случајност. Зато се сви добијени модели ОЗП обавезно евалуирају, по могућности статистичким тестовима, како би се сазнало да ли је откривена законитост заиста законитост или је производ погрешно постављеног експеримента или лоше изабраног узорка података.

ОЗП користи скуп алгоритама који захтевају одређена знања од корисника:

\begin{itemize}
  \item разумевање предуслова за рад алгоритма (захтевани тип података, дозвољену количину недостајућих података, захтевану припрему података),
  \item разумевање параметара алгоритма и делимично самог алгоритма,
  \item способност тумачења резултата добијених алгоритмом.
\end{itemize}

\subsection{CRISP-DM методологија}
Да би се ефикасно водили пројекти ОЗП постоје разне методологије. Једна од најпознатијих је \textbf{CRISP-DM} (међуиндустријски стандардни процес за откривање законитости у подацима). Састоји се из следећих фаза:

\begin{itemize}
  \item Разумевање пословног проблема,
  \item Разумевање података,
  \item Припрема података,
  \item Моделовање решења,
  \item Евалуација модела,
  \item Примена модела.
\end{itemize}
Методологија CRISP-DM је итеративног карактера са циљем да се дође до система ОЗП који је ефикасан и ефективан.

\subsubsection{Разумевање пословног проблема}
Прва фаза има за циљ да аналитичар добро разуме пословни проблем, да схвати који је циљ који треба да оствари будућим системом ОЗП. У овој фази развојни тим треба од стране доменских експерата да сазна захтеве које систем ОЗП треба да испуни. Ово је фаза у којој би требало да подстакне учење развојног тима, али и доменских експерата који често нису свесни домета које систем ОЗП може да оствари.

У овој фази се дефинишу циљеви и хипотезе истраживања, а такође се описују основни појмови из области у којој се ради анализа.

\subsubsection{Разумевање података}
Фаза у којој се аналитичар упознаје са подацима са којима треба да ради и да према томе одреди који му алгоритми, методе, архитектуре стоје на располагању. У овој фази се добијају сазнања о квалитету, формату и употребљивости података за анализу.

Аналитичар треба да обрати пажњу на:

\begin{itemize}
  \item \textbf{Тип података}:
  \begin{itemize}
    \item \textbf{Категорички} — ненумерички, без могућности успостављања поретка међу категоријама (нпр. за атрибут тип писаљке, категорије могу бити оловка, хемијска, фломастер).
    \item \textbf{Ординални} — категорички подаци у којима међу категоријама може да се успостави одређени поредак (нпр. вруће, топло, свеже, хладно).
    \item \textbf{Нумерички} — поред особина једнаког растојања међу бројевима имају и особину смислене нулте вредности.
  \end{itemize}
  \item \textbf{Смер података}:
  \begin{itemize}
    \item \textbf{Улазни} (познати и као дескриптори) — служе за прављење модела за предвиђање излазних атрибута, или се користе самостално.
    \item \textbf{Излазни} (познати као предиктори) — подаци које је потребно предвидети на основу улазних података.
  \end{itemize}
  \item \textbf{Корелација међу подацима} — мери се најчешће коефицијентом корелације, са вредностима у распону од -1 до 1, где 0 означава одсуство корелације, 1 апсолутну корелисаност, а -1 апсолутно негативну корелисаност. Корелација се испитује између парова улазних атрибута, излазних атрибута, као и улазних и излазних атрибута.
  \item \textbf{Дистрибуција (расподела) података} — колико су подаци учестали за конкретне вредности. Најчешће алгоритми претпостављају нормалну расподелу података (у којој се аритметичка средина најчешће појављује), мада постоје и многе друге могуће расподеле података (нпр. униформа у којој свака вредност може да се појављује са истом учесталошћу), а аналитичар би требало да буде упознат са овим расподелама.
  \item \textbf{Недостајуће вредности} — треба одлучити да ли редове са недостајућим подацима избацити из анализе или применити неку технику попуњавања (аритметичка средина за нумеричке, модус за категоричке).
  \item \textbf{Екстремне вредности} — аналитичар треба да донесе одлуку да ли да задржи аномаличне вредности у подацима или да их избаци и прогласи екстремима.
\end{itemize}
Аналитичар у фази разумевања података користи графике, али може да користи и алгоритме машинског учења како би стекао увиде у квалитет података и прелиминарну могућност аналитичке обраде истих.

\subsubsection{Припрема података}
Основни задатак припреме података јесте да се подаци за анализу сместе у табеларни облик (где колоне представљају атрибуте, а редови случајеве) или неки други погодан облик над којим је могуће спровести алгоритме ОЗП и открити законитости. Активност припреме података може да се схвати као функција која податке из неструктурираног облика претвара у структуриран облик за алгоритме, тј. облик неопходан за успешан рад алгоритама.

Сваки алгоритам захтева податке у одређеном формату и поседује претпоставке са каквим особинама података треба да ради. Алгоритми, осим ако то софтверски није подржано, не проверавају да ли је посао припреме података добро урађен — управо је посао аналитичара да податке доведе у исправан облик. Ова фаза, заједно са разумевањем података и пословног проблема, одузима знатан део времена у пројекту развоја система ОЗП.

Различито припремљени подаци могу да произведу различите излазе из модела, тако да треба наћи праву меру између степена припреме података (чиме се најчешће губи одређени квалитет података) и задржаног квалитета података који гарантују проналажење одређених законитости.

Технике за припрему података се грубо могу поделити у две категорије:

\begin{itemize}
  \item \textbf{Рад са редовима (случајевима)}:
  \begin{itemize}
    \item Селекција — креирање упита који извлачи одређени подскуп занимљив за анализу,
    \item Узорковање — узимање узорка из скупа података,
    \item Балансирање излазног атрибута — кад алгоритам тражи једнак или сличан број категорија по излазном атрибуту,
    \item Агрегирање — рачунање одређене сумарне статистике као што је нпр. центроид код алгоритма К-средина.
  \end{itemize}
  \item \textbf{Рад са колонама (атрибутима)}:
  \begin{itemize}
    \item Дефинисање типа и смера атрибута (могуће је исте податке посматрати или као категоричке или као ординалне или као нумеричке. Нпр. оцена на испиту може бити све три категорије, при чему различити алгоритми захтевају различити тип података),
    \item Филтрирање — избацивање нежељених атрибута,
    \item Извођење нових атрибута — на основу постојећих атрибута. Овај процес је изузетно значајан јер може да произведе атрибуте на основу којих алгоритам може да открије корисније законитости. Код дефинисања ових атрибута помоћ експерата је од пресудног значаја.
    \item Дискретизација — трансформисање нумеричких података у ординалне, тј. у распоне података или корпе. Овај процес као и извођење нових атрибута може помоћи у откривању значајнијих законитости.
  \end{itemize}
\end{itemize}

\subsubsection{Моделовање решења}
Када су схваћени пословни проблеми и подаци, и када су подаци припремљени, могуће је покренути алгоритме ОЗП. Процес моделовања решења представља централни део у изради пројекта ОЗП и најчешће се на њега мисли када се каже ОЗП. Ипак, овај део узима најмањи део времена у једном пројекту. Учешће људског фактора односи се на избор једног или више алгоритама и подешавање параметара изабраних алгоритама. У оквиру ове фазе битно је поставити и експеримент (нпр. унакрсна валидација, раздвајање података на скуп података за учење и проверу) за евалуацију модела ОЗП и развити такав модел који постиже математички одговарајућа мерила квалитета. Обично се овај процес евалуације означава као \textbf{верификација} — провера да ли машински научен модел ОЗП математички задовољава прагове квалитета.

\subsubsection{Евалуација система ОЗП}
Евалуација решења представља активност која треба да испита валидност (употребљивост, исправност) добијеног решења са стране корисника будућег система ОЗП. Математичка евалуација развијеног модела је део фазе моделовање решења. У овој фази се ради \textbf{валидација система ОЗП} — тј. да ли систем ОЗП може да унапреди у реалној пракси процес одлучивања и тиме оствари циљеве који су постављени у фази разумевања пословног проблема.

\subsubsection{Примена модела}
У процесу примене решења може да се сазна колико је систем ОЗП заиста употребљив. У овој фази се решење примењује над реалним окружењем, с тим да и даље постоји могућност да решење није у потпуности задовољавајуће, те се током рада адаптира и модификује. Током реалне примене система ОЗП константно се прате перформансе система, а систем се са времена на време унапређује, модификује или одбацује.

\subsection{Задаци ОЗП}
Задаци ОЗП представљају области примена система ОЗП. Сваки систем ОЗП може да има за циљ испуњавање једног или више задатака:

\begin{itemize}
  \item \textbf{Редукција} — свођење података на димензије које могу да се анализирају на ефикасан начин.
  \item \textbf{Регресија} — учење модела ОЗП који моделује везу између улазних атрибута и излазног атрибута који је нумеричког типа.
  \item \textbf{Класификација} — учење модела ОЗП који моделује везу између улазних атрибута и излазног атрибута који је категоричког или ординалног типа.
  \item \textbf{Кластеровање} — груписање случајева према сличности у одређени број група које нису унапред познате. Број група може да се одреди унапред, али и може да се врши претрага оптималног броја група, тј. кластера.
  \item \textbf{Асоцијативна и класификациона правила} — откривање ако-тада правила у структурама табелама и трансакцијама.
  \item \textbf{Предвиђање} — решавање свих претходних задатака у неком будућем временском периоду, нпр. за наредни дан, за наредну недељу, месец, годину.
\end{itemize}
Број задатака овим списком није исцрпљен, могући су разни хибридни задаци који комбинују наведене задатке са просторном и временском димензијом.

\subsubsection{Редукција/Оптимизација избора атрибута и случајева}
Основни циљ редукције је смањивање димензије проблема на онај облик који је повољан за анализу и који успоставља компромис између брзине решавања проблема и квалитета података за анализу. Може бити:

\begin{itemize}
  \item \textbf{Редукција атрибута} или селекција атрибута је оптимизациони проблем чији је циљ изабрати онај скуп атрибута који налази оптималан однос између жељених мера квалитета и количине неопходних атрибута. У данашње време процеси избора оптималног скупа атрибута су постали изузетно успешни и омогућавају да се простор оригиналних атрибута знатно смањи или замени новим, тзв. „скривеним“ слојем атрибута који одлично репрезентује оригиналне атрибуте.
  \item \textbf{Редукција случајева (редова)} је процес узорковања података, са циљем да се нађе узорак који веродостојно репрезентује претпостављену популацију. Овај задатак је исто изузетно сложен јер некад и не постоји довољно редова за анализу, па је потребно генерисати вештачки нове редове не би ли се дошло до одређеног броја случајева.
\end{itemize}

\subsubsection{Процена (регресија)}
Процена је задатак који има за циљ да открије модел који постоји између улазних атрибута и излазног атрибута нумеричког типа. Решава се коришћењем линеарних модела (линеарна регресија) или сложенијих алгоритама (вештачке неуронске мреже) које представљају универзалне апроксиматоре. Линеарни модели су по правилу апроксимација неког реалног нелинеарног модела, али се широко користе јер су једноставнији и разумљивији.

\subsubsection{Класификација}
Класификациони алгоритми уче моделе који праве функцију између улазних атрибута са излазним категоричким атрибутом. Излазни категорички атрибут је познат као класни атрибут. Задатак класификације је да открије законитост по којој се одређени објекат (ред у табели) сврстава у одређену класу.

Ако излазни атрибут има само две категорије (нпр. да/не) онда се ради о бинарној класификацији. У осталим случајевима ради се о вишеструкој класификацији.

Ако је потребно предвидети више излазних атрибута истовремено тада се користе алгоритми класификације за више излазних атрибута (више етикета).

\subsubsection{Кластеровање}
Док класификација служи да открије модел по коме се објекти распоређују у унапред познате класе (надгледано учење), кластеровање треба да пронађе законитост на основу које се објекти групишу у класе које нису познате на почетку процеса. Кластер је класа која није позната на почетку процеса тражења законитости, већ треба да се открије самим процесом кластеровања, да се протумачи и верификује. Због недостатка излазног атрибута, као у случају класификације, спада у класу алгоритама ненадгледаног учења. Ненадгледано учење значи да се у подацима не тражи повезаност између улазних атрибута и неког излазног атрибута, већ се у недостатку тог атрибута, заправо случајеви групишу по сличности.

Кластеровање је процес откривања група таквих да су објекти унутар кластера међусобно слични, а да су притом прилично различити од објеката других кластера. (Ово је дефиниција која је карактеристична за алгоритам К-средина, док други алгоритми кластеровања могу да претпостављају и различите дефиниције кластера.)

Кластеровање се спроводи из следећих разлога:

\begin{itemize}
  \item када нису унапред познате класе за разврставање објеката,
  \item пружа више информација о природи груписања објеката, тј. припадност кластера може даље да се користи као ново сазнање у процесу ОЗП. То чак може да се тумачи као нека „скривена“ особина.
  \item омогућава да се редукује број случајева који се анализирају тако што се уместо случајева који се налазе у кластеру узима само њихов представник. Тако је могуће заправо процес редукције случајева урадити уз помоћ алгоритама кластеровања при чему је потребно одредити број кластера који се жели имати. Представници ових кластера могу после да постану редуковани случајеви,
  \item омогућава организацијама да боље разумеју специфичности својих корисника.
\end{itemize}
Након развоја модела кластеровања, сваком случају може да се додели одређени излазни атрибут — припадност кластеру — чиме се отвара пут за коришћење техника класификације које могу да помогну да се открију законитости по којој су објекти кластеровани.

\subsubsection{Асоцијативна и класификациона правила}
Код класичне статистичке анализе аналитичар или ДО сами претпостављају правило (тј. хипотезу), док код ОЗП аналитичар задаје само параметре претраге за одређеним правилима, а алгоритам самостално открива законитости.

Асоцијативна правила су облика \textbf{АКО узрок ТАДА последица} и према томе су лако схватљива и применљива за моделовање знања. Битна су за разумевање и унапређење пословног процеса — омогућавају да се уоче неке законитости које се иначе тешко уочавају.

Класични проблеми који се решавају асоцијативним правилима су:

\begin{itemize}
  \item \textbf{Анализа потрошачке корпе} која може да одговори на питања који производи се купују заједно, или да ли куповина одређеног производа утиче на куповину другог производа.
  \item \textbf{Системи препоруке} — на основу куповине одређених производа препоручити потенцијално занимљив производ кориснику на основу сличних образаца куповине других корисника или особина које сам корисник поседује.
\end{itemize}

\subsubsection{Предвиђање}
Сваки претходно објашњени задатак истовремено може да служи и за предвиђање, поред алгоритама временских серија који природно у себи укључују и временску димензију. Постоје алгоритми за редукцију, процену, класификацију, кластеровање и откривање асоцијативних правила са подацима са временском димензијом али је могуће и инжењерингом података развити моделе за предвиђање будућих стања посматраног система.

\section{Алгоритми за откривање законитости у подацима}
У наставку ће бити представљено неколико алгоритама за откривање законитости у подацима, груписаних по следећим типовима:

\begin{itemize}
  \item \textbf{Алгоритми за компресију и откривање скривених фактора} — служе да се нађе компресована репрезентација оригиналних података таква да смањи величину улазног простора што више, а да задржи квалитет декомпресије у оригинални простор података.
  \item \textbf{Стабла одлучивања} — решавају задатак класификације и регресије, уче модел у облику дрвета одлучивања.
  \item \textbf{Асоцијативна и класификациона правила} — откривају правила облика Ако-Тада сходно постављеним праговима подршке и поверења (ови термини ће бити објашњени у наставку).
  \item \textbf{Алгоритми за кластеровање} — откривају груписања у подацима.
  \item \textbf{Регресиони алгоритми} — уче модел зависности између улазних и излазног податка нумеричког типа.
  \item \textbf{Вештачке неуронске мреже} — представљају универзални апроксиматор (може до перфекције да се научи над скупом података за учење) за решавање задатака класификације и процене.
\end{itemize}

\subsection{Редукција}
Постоји мноштво алгоритама који се користе за редукцију. Овде ћемо се концентрисати на редукцију атрибута, јер се алгоритми који се овим баве користе стандардно у процесима ОЗП.

Неки од најпознатијих алгоритама редукције који врше компресију улазних података у тзв. скривени (латентни) слој су:

\begin{itemize}
  \item \textbf{Анализа главних компоненти} — оригиналан скуп атрибута трансформише у скуп главних компоненти таквих да када се оригинални подаци пројектују на главне компоненте, међу њима не постоји зависност (корелација је једнака 0), док је варијабилитет појачан на првим компонентама. Открива се онолико компоненти колико има почетних атрибута, али се најчешће користило само 20\% компоненти да опише 80\% варијансе која се налази у подацима. У данашње време модерни алгоритми могу у неким случајевима да открију 1\% компоненти које описују 99\% варијансе која постоји у подацима.
  \item \textbf{Факторска анализа} — слично анализи главних компоненти ради екстракцију скривених фактора, али те факторе тражи над заједничком варијансом, а не укупном као код анализе главних компоненти. Ова наизглед мала разлика је суштинска, јер скривени фактори над заједничком варијансом омогућавају интерпретабилност ових фактора. У друштвеним наукама се преферира факторска анализа над анализом главних компонената чије откривене компоненте не морају нужно да имају значење у реалности.
\end{itemize}

\subsection{Асоцијативна и класификациона правила}
Асоцијативна правила представљају један од основних облика законитости какве се откривају код ОЗП — облик АКО-ТАДА. Алгоритми ОЗП откривају законитости у облику АКО (категорија, или конјункција категорија) ТАДА (категорија или конјункција категорија). Њихов циљ је да нађу скуп правила који, за задате параметре, откривају све постојеће АКО-ТАДА релације.

Асоцијативна правила представљају једну од првих техника које су се користиле код ОЗП. У наставку ће бити представљен један од најпознатијих алгоритама за откривање асоцијативних правила, Apriori.

\subsubsection{Разумевање података}
Алгоритам ради искључиво са категоричким (ненумеричким) подацима, тако да је корисно све податке припремити тако да буду ненумеричког типа. Уколико у скупу података постоје нумерички подаци, потребно је те податке претворити у категорички тип података.

\subsubsection{Трансформација података}
У случају да су атрибути нумеричког типа, потребно је дискретизовати податке, тј. нумеричке податке претворити у дискретне вредности. Неке од основних техника:

\begin{itemize}
  \item \textbf{Аритметичка средина ± n × стандардна девијација} — могуће је направити 3 интервала, 5 интервала и слично. Тако нпр. 3 интервала овом методом имаће следеће границе: (-∞, средина − девијација], (средина − девијација, средина + девијација], (средина + девијација, +∞).
  \item \textbf{Једнака фреквенција} — могуће је направити предефинисани број категорија тако да у сваком интервалу буде подједнак број елемената.
  \item \textbf{Кориснички дефинисани интервали} — границе одређује експерт на основу доменског знања.
\end{itemize}
\textbf{Пример:} атрибут потрошња минута разговора дневно за седам корисника, са вредностима: 50, 60, 65, 70, 80, 120 и 150. Циљ је дискретизација у три категорије: ниска, средња и висока.

\begin{itemize}
  \item \textbf{Метод средине и девијације:} аритметичка средина ≈ 85, стандардна девијација ≈ 35. Границе: (−∞, 50], (50, 120], (120, +∞). Вредност 50 → ниска; 60, 65, 70, 80, 120 → средња; 150 → висока.
  \item \textbf{Метод једнаке фреквенције:} сортирати растуће и поделити на три групе. Вредности 50 и 60 → ниска; 65 и 70 → средња; 80, 120 и 150 → висока.
  \item \textbf{Кориснички дефинисани интервали:} нпр. ≤ 60 ниска; (60, 100] средња; > 100 висока. Вредности 50 и 60 → ниска; 65, 70, 80 → средња; 120 и 150 → висока.
\end{itemize}

\subsubsection{Моделовање решења}
Да би алгоритам Apriori могао да почне своју претрагу потребно је дефинисати два параметра:

\begin{itemize}
  \item \textbf{Подршка} представља вероватноћу појављивања одређене категорије у подацима. Подршка може да се односи како на АКО део правила, тако на ТАДА део правила, али и на цело АКО-ТАДА правило. Говори о томе колико је правило или део правила заступљен у подацима. Обично се користи подршка АКО дела правила као почетни параметар, али у зависности од софтвера који се користи то може бити и другачије.
  \item \textbf{Поверење} представља условну вероватноћу АКО-ТАДА правила, тј. вероватноћу појављивања одређеног ТАДА дела правила у подскупу података АКО дела правила.
\end{itemize}
\textbf{Пример 1.} Анализира се 1000 трансакција куповина производа. Подршка куповине чоколаде $P(\textrm{чоколада})$ је 5\%, што значи да се у 50 куповина појављивала чоколада. Поверење куповине негазираног сока у куповинама чоколаде $P(\textrm{негазирани сок} \mid \textrm{чоколада})$ је 60\%, што значи да се у 50 куповина чоколаде негазирани сок појавио 30 пута.

Кораци алгоритма:

\begin{itemize}
  \item Претражити над свим категоријама (подскуповима) (почевши од скупова са једним елементом) инкрементално увећавајући скуп за по један елемент док год су подскупови заступљени задатом подршком. Скупови података (категорија) који задовољавају задату подршку се називају и фреквентни (учестали) скупови података.
  \item Унутар свих учесталих скупова података проверити да ли је између подскупова задовољено задато поверење. Сачувати она правила над којима је задовољено поверење.
\end{itemize}
Сам процес је рачунски врло захтеван јер се претражују све могуће комбинације подскупова података. Убрзавање поступка омогућава следеће правило:

\textit{Подршка било ког учесталог подскупа података мора имати подршку мању или једнаку учесталом скупу у коме је садржан подскуп.}

Ово омогућава да се простор претраге смањи јер нема сврхе претраживати оне подскупове чији скупови нису имали одговарајућу подршку.

\textbf{Пример 2.} Посматрамо 1000 трансакција куповина у супермаркету. Параметри: подршка = 0.3, поверење = 0.6.

\begin{itemize}
  \item У првој итерацији претражују се сви скупови са једним производом. Откривене подршке: $P(\textrm{пециво}) = 30\%$, $P(\textrm{млечни производи}) = 25\%$, $P(\textrm{сок}) = 33\%$.
  \item У следећој итерацији се претражују подскупови са два производа који имају степен подршке најмање 30\%. Види се да $P(\textrm{пециво, млечни производи}) = 18\%$, $P(\textrm{пециво, сок}) = 20\%$, $P(\textrm{млечни производи, сок}) = 15\%$ не могу да остваре жељени ниво подршке те само правило (пециво, сок) може имати тражену подршку у даљим итерацијама.
  \item Проверава се поверење: $P(\textrm{пециво} \mid \textrm{сок}) = 55\%$, $P(\textrm{сок} \mid \textrm{пециво}) = 65\%$. Откривена је законитост $P(\textrm{сок} \mid \textrm{пециво})$ која говори да се од свих куповина пецива сок купује у 65\% случајева. $P(\textrm{АКО пециво ТАДА сок}) = 65\%$.
\end{itemize}

\subsubsection{Евалуација модела}
Код асоцијативних правила је изузетно битно проверити ваљаност датог правила процесом верификације (мерама квалитета) и валидације (потврдом од стране експерата). Да би се проценио квалитет користе се три основне мере:

\textbf{Подршка} (Support) — мери колико се често правило појављује у скупу података.

\[
\textrm{подршка}(X \Rightarrow Y)=\frac{\vert \{t \in T : X \cup Y \subseteq t\}\vert }{\vert T\vert }
\]

где је T скуп свих трансакција, t једна трансакција, \(X \cup Y\) скуп ставки у правилу. Узима вредности између 0 и 1. Већа вредност значи да је правило чешће у подацима. Правила са веома малом подршком се често одбацују.

\textbf{Поверење} (Confidence) — мери колико често се ТАДА део правила Y појављује у трансакцијама које већ садрже АКО део X.

\[
\textrm{поверење}(X \Rightarrow Y)=\frac{\textrm{подршка}(X \cup Y)}{\textrm{подршка}(X)}
\]

Узима вредности између 0 и 1. Високо поверење значи да правило има добру предиктивну моћ. Међутим, високо поверење не мора нужно значити и јаку повезаност између X и Y.

\textbf{Надмашивач} (Lift) — мери колико је појава X и Y заједно већа (или мања) од очекиване вредности у случају њихове независности.

\[
\textrm{подизач}(X \Rightarrow Y)=\frac{\textrm{поверење}(X \Rightarrow Y)}{\textrm{подршка}(Y)}=\frac{\textrm{подршка}(X \cup Y)}{\textrm{подршка}(X) \cdot \textrm{подршка}(Y)}
\]

\begin{itemize}
  \item Надмашивач = 1 → не постоји веза између АКО и ТАДА дела правила.
  \item Надмашивач > 1 → постоји позитивна повезаност (АКО део повећава вероватноћу ТАДА дела).
  \item Надмашивач < 1 → АКО део правила утиче негативно на ТАДА део.
\end{itemize}
\textbf{Пример:} имамо 100 трансакција. У 30 се појављује X, у 40 се појављује Y, у 20 се заједно појављују X и Y.

\begin{itemize}
  \item Подршка = $20 / 100 = 0.20$
  \item Поверење = $20 / 30 \approx 0.67$
  \item Надмашивач = $0.67 / 0.40 \approx 1.67$
\end{itemize}
Резултати показују да се правило појављује у 20\% свих трансакција, у 67\% случајева када се појави X појављује и Y, а повезаност је позитивна и значајна (надмашивач > 1).

\subsubsection{Примена модела}
Асоцијативна правила могу да се користе у реалним ситуацијама као системи препоруке — да се корисницима јављају препоруке о евентуално занимљивим производима на основу до сада прегледаних. Користе се и за прављење база правила које се покрећу у реалном времену. Седамдесетих година прошлог века су експертски дефинисана асоцијативна правила коришћена за развој експертних система, а данас се често користе у системима извештавања у реалном времену (нпр. препоруке возачима за безбедну вожњу на аутопуту).

\subsection{Алгоритми за кластеровање}
Кластер алгоритми служе да случајеве у скупу података групишу у подскупове који су хомогенији од почетног скупа. Постоје различите фамилије алгоритама кластеровања:

\begin{itemize}
  \item \textbf{Засновани на представницима} — модел кластера се представља преко једног представника који је најчешће централни објекат кластера (нпр. аритметичка средина - центроид, или најцентралнији реални случај у кластеру - медоид). Најпознатији представник је алгоритам К-средина (k-means).
  \item \textbf{Хијерархијски} — сличности међу кластерима представљају у угњежденој хијерархијској структури тако да кластер на врху стабла садржи цео скуп података, а најнижи чворови представљају саме елементе скупа. Најпознатије врсте су агломеративни (полазе од претпоставке да су сви елементи посебан кластер и итеративно их уједињују у један кластер) и дивизиони (полазе од претпоставке да су сви елементи у истом кластеру и итеративно их деле до појединачних случајева). И код дивизионог и хијерархијског кластеровања, стабло представља процес угњеждавања/рашчлањивања објеката од целог скупа до појединих случајева или обрнуто.
  \item \textbf{Засновани на густини} — кластер се дефинише као подскуп који у одређеној околини има довољан број елемената. Алгоритам повезује све елементе скупа који имају сличну густину у један кластер.
\end{itemize}
Наведене три групе кластера имају сасвим другачији поглед на оно што јесу кластери, те се у реалној пракси често комбинују или користи она фамилија алгоритама која одговара проблему који се решава.

Сви кластер алгоритми откривају законитости груписања. Законитост код кластер алгоритма представља правило по коме се случај додељује кластеру.

Битно је напоменути да не постоји најбоље кластеровање, већ је најприхватљивији кластер модел онај који за доносиоца одлука има највише смисла и на основу кога могу да се доносе квалитетне одлуке.

\subsection{Алгоритам К-средина}

\subsubsection{Разумевање пословног проблема}
Једна од најчешћих примена кластеровања је сегментација корисника, како би се боље сазнале специфичне особине појединих група. Различите организације имају различите стратегије — неке праве један исти производ за све кориснике, а неке желе да омогуће да сваки клијент има себи својствен производ. И један и други сценарио делују прилично утопистички те се већина организација одлучује за управљање одређеног броја сегмената преко којих даље може да се боље прилагоди својим корисницима, кроз посебне пакете услуга и производа. Кластеровање може математички да помогне да се открију законитости груписања корисника.

Кластеровање је вишеатрибутивно груписање објеката (за разлику од сортирања које не може истовремено да сортира по више атрибута, већ се то ради лексикографски, тј. по одређеном редоследу).

Кластеровање је вероватно најзаступљенија фамилија алгоритама ОЗП, јер поред свог задатка ОЗП помаже и решавању других задатака ОЗП. Нпр. кластеровање се широко користи и у анализи слика (рендгенски снимци, сателитски снимци) ради идентификације занимљивих сегмената и откривање законитости таквих група.

\subsubsection{Разумевање података}
Подразумевано, алгоритам К-средина ради са нумеричким подацима. Мера сличности коју користи је еуклидско одстојање (корен збира квадратних одстојања) да би се одредило одстојање између сваког случаја и представника група. Због особине да користи еуклидско одстојање, алгоритам је добар за проналажење кластера сферног облика, тако да неће радити ефикасно са кластеровањем нпр. слика у којима треба да се препозна несферичан облик. (Ипак, различите имплементације омогућавају мерење сличности и помоћу других мера сличности, те проблем налажења несферичних облика може да се предупреди.)

Алгоритам је изузетно осетљив на:

\begin{itemize}
  \item \textbf{Атрибуте различитих величина скала} — ако је један атрибут изражен у $[10^3, 10^4]$, а други у децималним $[10^{-1}, 10^{-2}]$, тада други атрибут неће уопште утицати на мерење сличности. Због овога је неопходно користити нормализацију, скалирање атрибута пре примене самог алгоритма.
  \item \textbf{Аномаличне (екстремне) вредности} — вредности које значајно одступају од осталих елемената скупа могу значајно да утичу на перформансе алгоритма к-средина. Нпр. ако би се као почетни представник (центроид) једног од к кластера изабрао случај који је аномаличан, могла би да се деси ситуација да тај кластер не би успео да привуче ниједан други елемент, чиме би такав кластер сам по себи остао бесмислен. Овај проблем се решава или прелиминарном редукцијом скупа података уклањањем свих аномаличних случајева, или „паметним“ одабиром почетних к представника.
\end{itemize}

\subsubsection{Трансформација података}
Алгоритам захтева да су сви подаци нумеричког типа. Ако постоји атрибут ненумеричког типа, могуће је извршити трансформацију категоричког у нумерички тип. То се најчешће ради бинарним кодирањем у којем свака категорија добије један бинарни запис.

Да би се сви нумерички подаци довели у исти распон, могуће је користити методе нормализације, скалирања:

\begin{itemize}
  \item Скалирање у односу на збир вредности одређеног атрибута (збир скалираних вредности ће бити 1).
  \item Скалирање у односу на највећу вредност одређеног атрибута (једна вредност ће бити 1, остале представљају процентуални удео).
  \item Скалирање у односу на аритметичку средину и стандардну девијацију (од почетне вредности се одузима аритметичка средина, разлика се дели стандардном девијацијом). Ова метода је врло корисна и за идентификацију аномаличних вредности једног атрибута — оне најчешће одступају две или више стандардних девијација у односу на аритметичку средину.
\end{itemize}

\subsubsection{Моделовање решења}
Алгоритам поседује параметар \textit{к} (број кластера) који треба да се зада пре него што алгоритам почне са радом. Пошто најчешће није очигледно колико кластера постоји у подацима (знамо да се тај број креће од 1 до укупног броја случајева у скупу података), често се алгоритам покреће више пута са различитим вредностима \textit{к} уз активно учешће доменског експерта. Док алгоритам к-средина може математички да одреди оптималан број кластера, доменски експерт је тај који одлучује да ли је решење прихватљиво, да ли су коришћени одговарајући атрибути за кластеровање, и слично.

Кораци алгоритма:

\begin{enumerate}
  \item Из скупа података насумично изабрати \textit{к} случајева и прогласити их центроидима (представницима).
  \item Све остале елементе скупа доделити најближим представницима тако што се сви случајеви доделе представнику од кога су најмање удаљени.
  \item Прерачунати к центроида као средње вредности случајева који су додељени одговарајућим кластерима.
  \item Понављати кораке 2 и 3 до конвергенције. (Ово значи да нема више прелажења случајева међу кластерима или да константно исти случајеви остају у истим кластерима.)
\end{enumerate}
Почетна фаза случајне иницијализације центроида је осетљива, те се често алгоритам покреће више пута за дефинисано \textit{к} да би се избегла ситуација да један од представника буде аномалија.

\subsubsection{Евалуација модела}
За верификацију модела кластеровања постоје бројне мере квалитета кластера.

\textbf{Индекс силуете/сенке} (Silhouette index) је једна од најпопуларнијих мера квалитета:

\[
s_{i}=\frac{b_{i}-w_{i}}{\max(b_{i},w_{i})}
\]

где је:

\begin{itemize}
  \item $b(i)$ (between, између) — просечна удаљеност случајева посматраног кластера од случајева у најближем суседном кластеру,
  \item $w(i)$ (within, унутар) — просечна међусобна удаљеност случајева унутар припадајућег кластера.
\end{itemize}
Индекс сенке проверава да ли је одређени елемент сличнији елементима унутар свог кластера или елементима најближег суседног кластера. Узима вредности у интервалу од -1 до 1, где 1 значи апсолутну припадност припадајућем кластеру, док -1 означава неприпадност. Срачунава се за сваки случај и на крају се рачуна средња вредност свих индекса сенки и тај сумарни број означава квалитет модела кластеровања. Пожељно је да ова вредност буде што ближа јединици. Вредности 0,5 и веће означавају да сви елементи у просеку више припадају свом кластеру него најближем суседном кластеру. У реалним применама вредности изнад 0,4 су исто прихватљиве.

\textbf{Инерција} представља збир квадрата свих унутар-кластерских растојања и она представља меру компактности појединих кластера. Дефинише се као збир квадрата растојања свих објеката од представника (центроида) кластера којем припадају. Нижа вредност значи да су објекти ближи својим центрима, односно да су кластери компактнији и хомогенији.

Важно је напоменути да се вредност инерције готово увек смањује са повећањем броја кластера. Због тога се инерција не може посматрати изолована, већ се најчешће користи у комбинацији са методама за избор броја кластера.

\textbf{Лакат метода} представља интуитиван приступ за одређивање оптималног броја кластера. Заснива се на анализи промене вредности инерције у зависности од броја кластера. Алгоритам кластеровања се покреће више пута за различите вредности броја кластера (нпр. од 2 до 10), при чему се за сваку вредност израчунава инерција. У почетку, повећање броја кластера доводи до значајног смањења инерције. Након одређеног броја кластера, даљи раст доноси све мања побољшања. Тачка на графику где долази до нагле промене у брзини опадања инерције подсећа на облик лакта, по чему је метода и добила име. Број кластера који одговара тој тачки сматра се оптималним.

\textbf{Валидација} је процес у коме корисник треба да потврди да пронађени кластери одговарају траженим сегментима и да на основу њих може да ефикасније одлучује. Корисник при процесу валидације прегледа центроиде и покушава да их протумачи. Ако су центроиди разумљиви и корисни доносиоцу одлука, кластер модел може да се прогласи прихватљивим. Наравно да корисника система занима и да ли је кластер модел математички прихватљив, али то се подразумева да је тако. Тек математички исправан кластер модел се пружа ДО на разматрање. ДО затим сваком кластеру може да додели одређене стратегије, тактике, производе, услуге, и/или да спроводи управљачке акције.

\subsubsection{Примена модела}
Модел к-средина се заправо састоји од \textit{к} центроида, тј. агрегираних (упросечених) представника својих кластера. Они представљају законитости модела кластеровања, тј. законитост по којој се случајеви расподељују у групе. У фази примене то значи да центроиди одређују над подацима над којима модел није учио који случајеве додељује ком кластеру.

Алгоритам К-средина спада у један од најкоришћенијих алгоритама код ОЗП. Често је такозвани основни алгоритам кластеровања, прва алтернатива код избора алгоритма за кластеровање објеката. Уколико не може да пронађе одговарајуће кластере, прелази се на друге фамилије кластер алгоритама.

Области примене алгоритма су широке: од кластеровања корисника, пацијената, радних процеса, текста, слика, К-средина се увек показује као истакнут алгоритам.

\subsection{Дрва одлучивања}
Дрва одлучивања производе законитости у форми дрва одлучивања. Свако дрво одлучивања састоји се од чворова и грана. Највиши чвор у хијерархији стабла се означава као корен или пањ, док се најнижи чворови у хијерархији стабла називају листови.

Свако стабло одлучивања се састоји од једног или више листова. Листови могу да се посматрају као асоцијативна правила. Међутим, за разлику од асоцијативних правила, стабла одлучивања повезују асоцијативна правила кроз хијерархијску форму дрва, тј. намећу хијерархијску структуру. Овим се добија механизам закључивања који једнозначно за одређене услове правила предлаже одређене одлуке. Код асоцијативних правила је могуће да се за исте услове покрене више правила истовремено.

Дрво одлучивања представља модел хијерархијски уређених асоцијативних правила који ефикасно решава проблеме класификације и регресије:

\begin{itemize}
  \item \textbf{Класификациони модел} је функција која за скуп различитих вредности улазних атрибута нумеричког и категоричког типа предлаже категорију излазног атрибута.
  \item \textbf{Регресиони модел} је функција која за одређени скуп различитих вредности атрибута нумеричког и категоричког типа предлаже нумеричку излазну вредност.
\end{itemize}
Постоји много алгоритама за дрва одлучивања. У наставку ће бити представљен један општи оквир алгоритама дрва одлучивања који решава проблеме класификације и може да ради и са нумеричким и категоричким типовима података.

\subsubsection{Разумевање пословног проблема}
Најпознатији проблеми класификације у пословној пракси су:

\begin{itemize}
  \item \textbf{Детекција одлазећих корисника} — проблем у коме треба да се развије модел који може да предвиди клијента који ће потенцијално да напусти компанију. Свака компанија тежи задовољству својих клијената и тежи да их задржи, јер се процењује да је мања инвестиција задржати постојећег клијента него стећи новог. Са становишта класификације ово представља проблем бинарне класификације. На основу историјских података учи се модел који препознаје клијенте који потенцијално напуштају организацију (Класа 1) или остају у организацији (Класа 0). За пословне кориснике је занимљиво да поред сигнала о напуштању сазнају и разлоге који су довели до одлуке о напуштању организације, а модели дрва одлучивања погодни су за решавање и оваквих проблема због своје особине разумљивости логике одлучивања алгоритма.
  \item \textbf{Детекција преваре} треба да открије аномалично понашање код корисника и да га спречи да изврши превару или да спречи да се над одређеним корисником изврши превара.
  \item \textbf{Процена кредитног ризика} је једна од активности које банке свакодневно спроводе и то скоро потпуно аутоматски. Изградњом класификационих модела процес процене кредитног ризика се у потпуности аутоматизује и тиме се знатно смањују трошкови људског доносиоца одлука.
\end{itemize}
Дрва одлучивања се намећу као један од неизбежних алгоритама код ОЗП из најмање два разлога:

\begin{itemize}
  \item разумљивости модела,
  \item могућности интеграције постојећих правила (експертског знања) и надоградње модела на основу историјских података.
\end{itemize}

\subsubsection{Разумевање података}
Алгоритам стабла одлучивања нема проблема са аномаличним подацима, а за недостајуће вредности стабла одлучивања имају способност да науче тзв. сурогат атрибуте за одлучивање. Ово значи да када наиђе један случај коме недостају одређене вредности, дрво одлучивања може у моделу да искористи и неки други атрибут којим може да се замени атрибут који недостаје.

Код стабла одлучивања проблем знају да буду категорички атрибути са великим бројем категорија, што доводи до превеликог броја грана које излазе из једног чвора и до ефекта тзв. широких, а плитких стабала, што зна да доведе до смањивања предиктивних перформанси модела.

Пошто алгоритам рачва категоричке атрибуте по свим гранама производећи плитка и широка дрва, а нумеричке атрибуте бинарно производећи уска и дубока дрва, тј. различито ради са нумеричким и категоричким атрибутима, треба посветити пажњу припреми података пре коришћења дрва одлучивања.

Још један од великих недостатака стабла одлучивања је рад са излазним атрибутом (атрибут који садржи класе) у ситуацији када је удео класа високо небалансиран. На пример, ако је једне класе 1\%, а друге класе 99\%, алгоритам би имао проблема да исправно изврши класификацију. Ипак, овакви проблеми се решавају померањем граница одлучивања за класификацију чиме стабла поново постају употребљива.

\subsubsection{Трансформација података}
Кључне трансформације:

\begin{itemize}
  \item \textbf{Груписање категорија} једног атрибута у мањи број категорија је процес који се извршава када постоји велики број категорија (нпр. списак позивних бројева телефона) и то зна да угрози тачност модела. Категорије могу да се групишу ручно (уз помоћ експерта) или аутоматски (нпр. путем алгоритма за асоцијативна правила, или путем хи-квадрат теста или слично).
  \item \textbf{Дискретизација нумеричких атрибута} представља превођење нумеричких у категоричке податке јер алгоритам стабла одлучивања првенствено ради са категоричким подацима, а и нумерички атрибути се обично увек рачвају код модела дрва одлучивања само бинарно.
  \item \textbf{Балансирање класа} се користи у ситуацијама велике неравнотеже између категорија излазног атрибута (тзв. класног атрибута). Могу да се користе технике узорковања:
  \begin{itemize}
    \item \textbf{Повећање броја елемената мањинске класе} — редови мањинске класе се поново узоркују.
    \item \textbf{Смањење броја елемената већинске класе} — скуп већинске класе се смањи на мањи број елемената.
  \end{itemize}
\end{itemize}

\subsubsection{Моделовање решења}
Алгоритам дрва одлучивања се састоји од неколико корака:

\begin{enumerate}
  \item За свако могуће рачвање (једнако свим расположивим рачвањима свих расположивих атрибута) рачуна се мера квалитета рачвања.
  \item Атрибут са најбољом мером квалитета се проглашава чвором.
  \item Скуп података се тада дели на подскупове који одговарају гранама најбољег чвора.
  \item Кораци 1--3 се понављају на свим подскуповима независно до испуњења критеријума заустављања раста стабла.
\end{enumerate}
Критеријуми заустављања раста стабла могу бити:

\begin{itemize}
  \item долазак до чистих чворова — чворова у којима излазни атрибут припада само једној класи (природан критеријум заустављања),
  \item када стабло достигне одређену дубину, што је критеријум који дефинише корисник.
  \item када чвор не садржи више од одређеног броја случајева, што је такође критеријум који дефинише корисник.
\end{itemize}
Нумерички атрибути се такође тестирају као кандидати за рачвање. Поступак проналаска оптималне тачке рачвања нумеричког атрибута:

\begin{enumerate}
  \item Нумерички атрибут се поређа у нерастући редослед и свакој вредности елемента се придружује класа којој елемент припада.
  \item Означавају се све бројне вредности у којима долази до промене између вредности класа. Доказано је да само у тим тачкама може доћи до оптималног рачвања.
  \item Рачуна се мера квалитета за свако рачвање и оптималним рачвањем се проглашава оно у коме је постигнута оптимална мера квалитета рачвања.
\end{enumerate}

\subsubsection{Евалуација модела}
Да би се оценио квалитет модела дрва одлучивања, потребно је проверити да ли класе које предвиђа модел дрва $\hat{y}$ одговарају стварној класи случајева $y$.

Основни начин за одређивање квалитета модела је да се изброји проценат поклапања предвиђених класа $\hat{y}$ и стварних класа $y$ (тачност класификације).

Треба водити рачуна да се тачност модела не мери искључиво над подацима над којима је учено дрво јер то доводи до нереалне процене квалитета рада модела над будућим новим случајевима, над којима дрво није учено.

Постоји неколико начина да се измери тачност класификације над подацима над којима стабло није учено:

\begin{itemize}
  \item \textbf{Раздвајање скупа података на податке за учење и податке за проверу} тако да се подаци за учење користе за учење модела, а подаци за проверу се користе да би се утврдило колико је модел тачан над подацима над којима није учио. Код ОЗП битно је да мере тачности класификације представљају процењену тачност класификације над новим/невиђеним подацима.
  \item \textbf{Дељење} скупа података може да се ради случајним путем или емпиријским (нпр. модел за кредитни ризик се развија над подацима из банке од 2000. до 2010, а тестира над подацима од 2011. до 2014).
  \item \textbf{Скуп података за учење, проверу и потврђивање} — поред скупа за проверу, може се увести и скуп за потврђивање (валидацију) над којим модел може бити пуштен само једном. Овакво тестирање представља неку додатну сигурност да тачност модела што више одговара ономе што може да се очекује када се модел пусти у реалну примену.
  \item \textbf{Унакрсна валидација} је математички најисправнији начин да се тачно утврди реална тачност рада модела над новим подацима над којима модел није учен и који се очекују у примени. Код унакрсне валидације скуп података за учење се дели на $n$ делова тако да се учи $n$ модела, сваки над $n - 1$ делом скупа података, а тестира над 1 делом. Сваки од $n$ модела се проверава сваки пут над различитим подскупом података и мери тачност класификације. Провера се увек обавља над другим подскупом и сваки пут се бележи тачност класификације. На крају се усредњена вредност ових тачности $n$ модела показала као најбољи начин да проценимо квалитет модела класификације.
\end{itemize}
Битно је напоменути да тачност класификације није добра мера квалитета модела у ситуацијама када постоји високо неуравнотежена расподела категорија излазног атрибута на скупу података. Када се дешава оваква ситуација треба се присетити основних статистичких мера за мерење квалитета модела — алфа и бета грешке.

\subsubsection{Матрица конфузије}
Код бинарне класификације (две класе 1 и 0) тачност модела се обично представља у \textbf{матрици конфузије}:

\begin{table}[h]
\centering
\caption{Матрица конфузије за бинарну класификацију}
\label{tab:matrica-konfuzije}
\begin{tabular}{|l|c|c|}
\hline
 & \textbf{Предвиђено: $\hat{y} = 1$} & \textbf{Предвиђено: $\hat{y} = 0$} \\
\hline
\textbf{Стварно: $y = 1$} & $a$ (ТП) & $b$ (НН) \\
\hline
\textbf{Стварно: $y = 0$} & $c$ (НП) & $d$ (ТН) \\
\hline
\end{tabular}
\end{table}
\begin{itemize}
  \item \textbf{ТП (тачно позитивни)} — број тачно предвиђених случајева класе 1
  \item \textbf{НН (нетачно негативни)} — број нетачно предвиђених случајева класе 1
  \item \textbf{НП (нетачно позитивни)} — број нетачно предвиђених случајева класе 0
  \item \textbf{ТН (тачно негативни)} — број тачно предвиђених случајева класе 0
\end{itemize}
Циљ сваке класификације је да има што више попуњена поља $a$ и $d$, а да смањи учешће поља $b$ и $c$. Тачност класификације се рачуна као $(a + d) / (a + b + c + d)$.

Често тај циљ није могуће остварити, већ је потребно дозволити једну грешку зарад смањења друге грешке. Две најбитније грешке су:

\begin{itemize}
  \item \textbf{Алфа грешка} или лажна узбуна је ситуација када се прогласи класом 1 оно што је заправо класа 0. Рачуна се као $c / (c + d)$ и представља грешку лажне узбуне (лажно се клијент проглашава да ће напустити компанију, или да је здрав пацијент болестан).
  \item \textbf{Бета грешка} је неспособност модела да предвиди класу од интереса, тј. класу 1. Рачуна се као $b / (a + b)$ и представља грешку пропуштене прилике (модел не може да предвиди да ће неко напустити компанију или да је неко болестан).
\end{itemize}
Алфа и бета грешка су међусобно негативно корелисане — желећи да смањимо једну грешку модела, често повећавамо другу.

Поред укупне тачности, у пракси се користе и додатне мере:

\textbf{Прецизност} — мери колико су предвиђања класе 1 тачна, односно који део елемената које је модел означио као класу 1 заиста припада тој класи:

\[
\textrm{прецизност}=\frac{\textrm{ТП}}{\textrm{ТП}+\textrm{НП}}=\frac{a}{a+c}
\]

Висока прецизност значи да модел ретко прави лажне узбуне (мала алфа грешка).

\textbf{Одзив} (осетљивост) — показује колико добро модел успева да препозна све елементе класе 1:

\[
\textrm{одзив}=\frac{\textrm{ТП}}{\textrm{ТП}+\textrm{НН}}=\frac{a}{a+b}
\]

Висок одзив значи да модел ретко превиђа позитивне случајеве (мала бета грешка).

\textbf{Ф1 мера} — представља хармонијску средину прецизности и одзива:

\[
F_{1}=2 \cdot \frac{\textrm{прецизност} \cdot \textrm{одзив}}{\textrm{прецизност}+\textrm{одзив}}=\frac{2a}{2a+b+c}
\]

Ф1 мера даје уравнотежену оцену модела када је потребно истовремено минимизовати и лажне узбуне и пропуштене случајеве.

\subsubsection{Примена модела}
Стабла одлучивања представљају једну од најпопуларнијих фамилија алгоритама за ОЗП. Снага им је у једноставности и могућности да производе интерпретабилне моделе. Велику примену имају у свим областима у којима се осим тачности класификације захтевају и разумљиви модели, тј. модели које доносилац одлуке може да прочита и разуме. Широку примену имају, између осталих, у областима медицине, маркетинга и тумачења пословних процеса.

\subsubsection{График надмашивања и кумулативних добити}
У пракси се често предвиђене вероватноће припадања класи од интереса (класа 1) за сваки случај скупа података користе интензивно. Да би се знало колико веровати вероватноћама $P(\hat{y} = 1)$ праве се две врсте кривих, и то крива надмашивања и крива кумулативне добити. У оба случаја се на апсциси (x оса) приказују предвиђања модела позитивне класе ($\hat{y} = 1$) за све случајеве у скупу података. Кључно је да се $P(\hat{y} = 1)$ поређају у опадајућем редоследу, тако да су веће вероватноће ближе координатном почетку, а мање вредности вероватноће су удаљеније од координатног почетка.

На ординати (y оса) се код криве кумулативне добити приказује стопа тачно предвиђених позитивних случајева, тј. удео тачно позитивних за посматрану вредност апсцисе у односу на укупан број позитивних. Ова крива би требало да буде права линија уколико је предвиђање перфектно, међутим у пракси се пре дешава да је крива конкавна и да у делу ближе координатном почетку очекујемо веће вредности на ординати, тако да може да се деси Парето принцип, нпр. да се у првих 20\% позитивно предвиђених вероватноћа нађе 80\% свих тачно предвиђених позитивних случајева.

Крива надмашивања је компатибилна са кривом кумулативних добити, тј. на ординати се приказује надмашивање, тј. колико је кумулативна добит већа од очекиване кумулативне добити за проценат посматране ситуације. За пример Парето 80\% откривених тачно позитивних у 20\% поређане популације позитивних, надмашивање је 4, тј. 4 пута је већа вероватноћа да се сусретне тачно позитивно предвиђање у 20\% популације са највећим вероватноћама $P(\hat{y} = 1)$. Знајући ово, организације често користе само случајеве са највишим вероватноћама за своје даље активности, знајући да тиме још више повећавају вероватноћу да процес одлучивања буде ефикасан.

\subsection{Линеарна регресија}
Линеарна регресија представља један од основних алгоритама за решавање проблема регресије, тј. предвиђања нумеричке вредности једног излазног атрибута на основу скупа улазних атрибута. За разлику од класификационих модела, код линеарне регресије излазна променљива је нумеричког типа.

Основна идеја линеарне регресије јесте да се пронађе линеарна функција која најбоље описује зависност између улазних атрибута и излазне променљиве. Модел линеарне регресије претпоставља да се излазна вредност може апроксимирати линеарном комбинацијом улазних атрибута.

\subsubsection{Разумевање пословног проблема}
Најчешћи проблеми регресије у пословној пракси су предвиђање прихода или продаје, процена вредности некретнина, предвиђање потрошње корисника, процена трајања одређених процеса, предвиђање метеоролошких параметара (падавине, температуре и сл.).

\subsubsection{Модел линеарне регресије}
У најједноставнијем случају проста линеарна регресија са једним улазним атрибутом има облик:

\[
y = \beta_{0} + \beta_{1} x
\]

где је $y$ излазна променљива, $x$ улазни атрибут, $\beta_0$ слободан члан, а $\beta_1$ коефицијент правца.

У општем случају са више атрибута модел линеарне регресије има облик:

\[
y = \beta_{0} + \beta_{1} x_{1} + \beta_{2} x_{2} + \ldots + \beta_{n} x_{n}
\]

Овај модел представља хиперраван у простору атрибута и параметри тог модела се уче током процеса машинског (статистичког) учења или се одређују аналитичким путем.

\subsubsection{Разумевање података}
Линеарни модел претпоставља:

\begin{itemize}
  \item линеарну везу између улазних атрибута и излазне променљиве,
  \item да атрибути нису јако међусобно корелисани (мултиколинеарност може довести до нестабилних и тешко интерпретабилних коефицијената атрибута),
  \item да подаци не садрже велики број екстремних вредности (могу имати несразмерно велики утицај на квалитет модела).
\end{itemize}
Уколико све ове претпоставке нису задовољене, перформансе модела могу бити значајно нарушене.

\subsubsection{Моделовање решења}
Коефицијенти линеарног модела могу се одредити на два основна начина: аналитичким путем или машинским учењем нумеричким методама оптимизације.

\textbf{Аналитичко решење} — проблем минимизације средње квадратне грешке решава се директно:

\[
\beta = (X^{T} X)^{-1} X^{T} y
\]

где је $X$ матрица улазних атрибута, $y$ вектор стварних вредности. Овај приступ даје тачно решење, али може бити рачунски захтеван код великих скупова података или када матрица $X^{T} X$ није инвертибилна.

\textbf{Градијентни спуст} је итеративни нумерички метод који се користи за минимизацију функције грешке:

\[
\beta = \beta - \alpha \nabla J(\beta)
\]

где је $\alpha$ брзина учења, а $\nabla J(\beta)$ градијент функције грешке.

Интуитивно, градијентни спуст се може посматрати као процес спуштања низ брдо. Замислимо да функција грешке представља површину у простору, где свака тачка на тој површини одговара одређеним вредностима коефицијената модела, а висина те тачке представља величину грешке. Циљ је да се пронађе најнижа тачка на тој површини, односно минимум функције грешке.

Алгоритам започиње у некој почетној тачки и у свакој итерацији одређује правац највећег опадања грешке (правац најстрмијег спуштања), који је одређен градијентом. Након тога се прави мали корак у том правцу, чиме се долази до тачке са мањом грешком. Поступак се понавља све док се не дође до тачке у којој даље спуштање више не доноси значајно побољшање.

Величина корака (брзина учења) има важну улогу: ако је корак превелик, може се прескочити минимум, док премали кораци доводе до спорог конвергирања.

\subsubsection{Евалуација модела}
\textbf{Средња квадратна грешка} мери просечну квадратну разлику између стварних и предвиђених вредности. Квадрирањем се додатно наглашавају веће грешке:

\[
\mathrm{MSE} = \frac{1}{n} \sum_{i=1}^{n} (y_{i} - \hat{y}_{i})^{2}
\]

\textbf{Корен средње квадратне грешке} има исту јединицу мере као и циљна променљива, што олакшава интерпретацију:

\[
\mathrm{RMSE} = \sqrt{\mathrm{MSE}}
\]

\textbf{Средња апсолутна грешка} рачуна просечну апсолутну разлику без квадрирања, те мање кажњава екстремне грешке:

\[
\mathrm{MAE} = \frac{1}{n} \sum_{i=1}^{n} \vert y_{i} - \hat{y}_{i} \vert
\]

\textbf{Коефицијент детерминације} показује колики део варијансе излазне променљиве модел успева да објасни:

\[
R^{2} = 1 - \frac{\sum_{i=1}^{n} (y_{i} - \hat{y}_{i})^{2}}{\sum_{i=1}^{n} (y_{i} - \bar{y})^{2}}
\]

Вредност се обично креће између 0 и 1. Већа вредност значи да модел боље објашњава варијансу. Вредност близу 1 значи врло добро уклапање. Вредност близу 0 означава да модел не успева значајно боље да објасни податке у односу на једноставан модел који увек предвиђа средњу вредност. У одређеним ситуацијама вредност може бити и негативна, тј. модел је гори од тривијалног приступа заснованог на просеку.

\textbf{Пример:} стварне вредности 10, 20, 30, 40. Предвиђене вредности 8, 25, 28, 50.

\begin{itemize}
  \item Грешке: $-2, +5, -2, +10$
  \item Квадрати грешака: $4, 25, 4, 100$. Збир $= 133$. $\mathrm{MSE} = 133 / 4 = 33{,}25$.
  \item $\mathrm{RMSE} \approx 5{,}77$.
  \item Апсолутне грешке: $2, 5, 2, 10$. Збир $= 19$. $\mathrm{MAE} = 19 / 4 = 4{,}75$.
  \item Средња вредност стварних података $= 25$. Збир квадрата одступања $= 500$. $R^2 = 1 - 133/500 \approx 0{,}73$ → модел објашњава око 73\% варијансе.
\end{itemize}

\subsubsection{Примена модела}
Након што је модел обучен и евалуиран, користи се за предвиђање вредности излазне променљиве на новим подацима. У пракси то значи да се за нове инстанце израчунава процењена вредност коришћењем научених коефицијената.

Примена подразумева да се нови подаци претходно трансформишу на исти начин као и подаци коришћени за тренирање — исти поступци скалирања, обраде недостајућих вредности и евентуалних трансформација. Уколико се ове трансформације не примене доследно, добијени резултати могу бити непоуздани.

Резултат је нумеричка вредност која представља процену циљне променљиве. На пример, у случају предвиђања продаје, модел може помоћи у планирању залиха, док у случају процене вредности некретнина може служити као подршка при одређивању цене.

\subsection{Логистичка регресија}
Логистичка регресија представља један од основних алгоритама за решавање проблема класификације. Инспирација за развој логистичке регресије је потекла из линеарне регресије у којој коефицијенти добијеног модела говоре о степену утицаја које одређени улазни атрибут има на категорију излазног атрибута. За разлику од линеарне регресије, која предвиђа нумеричку вредност излазног атрибута, логистичка регресија предвиђа вероватноћу одређене категорије излазног атрибута.

\subsubsection{Разумевање пословног проблема}
Логистичка регресија се користи у ситуацијама када се испуњава задатак класификације, као што је: да ли ће клијент напустити компанију, да ли треба одобрити кредит, да ли је трансакција превара, да ли је пацијент болестан. Модели дрва одлучивања и логистичке регресије решавају исти задатак с тим да су научени модели различити. Док код дрвета одлучивања листови одређују модел класификације, дотле код логистичке регресије то чине коефицијенти атрибута модела.

\subsubsection{Модел логистичке регресије}
За разлику од линеарне регресије, логистичка регресија не моделује директно излазну променљиву, већ вероватноћу припадности класи од интереса, нпр. 1. Да би се обезбедило да излаз модела буде између 0 и 1, користи се \textbf{логистичка (сигмоидна) функција}:

\[
P(y = 1 \mid x) = \frac{1}{1 + e^{-(\beta_{0} + \beta_{1} x_{1} + \ldots + \beta_{n} x_{n})}}
\]

Линеарна комбинација атрибута трансформише се сигмоидном функцијом у вероватноћу:

\begin{itemize}
  \item Ако линеарни део модела даје вредност 0, логистичка функција враћа вероватноћу 0,5 — модел је неодлучан.
  \item Ако даје позитивну вредност 2, вероватноћа је око 0,88 — модел процењује око 88\% вероватноће да пример припада класи 1.
  \item Ако даје вредност 4, вероватноћа је око 0,98 — модел је веома сигуран да пример припада класи 1.
  \item Ако даје негативну вредност $-2$, вероватноћа је око 0,12 — пример би био класификован као класа 0.
  \item Ако даје $-4$, вероватноћа је око 0,02 — модел практично са великом сигурношћу додељује пример класи 0.
\end{itemize}
Логистичка регресија се може посматрати као проширење линеарне регресије, где се резултат линеарне функције провлачи кроз функцију која га ограничава на интервал између 0 и 1. На тај начин модел не даје произвољне нумеричке вредности, већ вероватноће које имају јасно значење у контексту класификације.

\subsubsection{Разумевање података}
Логистичка регресија претпоставља приближно линеарну везу између атрибута и логаритма шансе (количник вероватноће да ће се остварити категорија од интереса и вероватноће да се неће остварити категорија од интереса), да атрибути нису јако међусобно корелисани, да подаци не садрже велики број екстремних вредности.

\subsubsection{Трансформација података}
Пре примене модела логистичке регресије често се врши:

\begin{itemize}
  \item \textbf{Скалирање атрибута} — да се сви нумерички атрибути доведу на приближно исту скалу.
  \item \textbf{Кодирање категоричких променљивих} је неопходно јер логистичка регресија ради само са нумеричким вредностима, док дрво одлучивања ради суштински само са категоричким подацима. Један од најчешћих приступа је да се свака категорија претвара у посебну бинарну променљиву.
  \item \textbf{Попуњавање недостајућих вредности} је изузетно битно код логистичке регресије јер се модел логистичке вероватноће учи на основу расподеле улазних вредности у односу на излазни атрибут. Недостајуће вредности могу озбиљно да поремете расподеле података, те самим тим модел који се учи није довољно репрезентативан.
  \item \textbf{Уклањање или третман изузетака} је истог значаја као и рад са попуњавањем недостајућих вредности. Аномалије је потребно уклонити из података да би научени параметри модела имали смисла и цео процес учења се спровео ефикасно.
\end{itemize}

\subsubsection{Моделовање решења}
За разлику од линеарне регресије, логистичка регресија не користи средњу квадратну грешку, већ функцију вероватноће — \textbf{логаритамску функцију губитка}:

\[
J(\beta) = -\frac{1}{n} \sum_{i=1}^{n} \left( y_{i} \log p_{i} + (1 - y_{i}) \log (1 - p_{i}) \right)
\]

где је $p_i$ вероватноћа коју модел додељује класи 1. Циљ модела је да максимизује вероватноћу тачних класификација, односно да минимизује логаритамску функцију губитка. Оптимизација се најчешће врши итеративним методама као што је градијентни спуст.

\subsubsection{Ограничења модела}
\begin{itemize}
  \item Модел претпоставља линеарну зависност између улазних атрибута и логаритма односа вероватноћа, што у пракси често није у потпуности задовољено.
  \item Осетљив је на мултиколинеарност, што може довести до нестабилних процена коефицијената.
  \item Не може да моделује сложене нелинеарне односе без додатних трансформација.
  \item У поређењу са сложенијим алгоритмима (ансамбли, неуронске мреже), може имати слабије перформансе у ситуацијама када су обрасци у подацима комплексни и нелинеарни.
\end{itemize}

\subsubsection{Евалуација модела}
Користе се већ описане мере засноване на матрици конфузије — тачност, прецизност, одзив и Ф1 мера.

\subsubsection{Примена модела}
За сваки нови пример модел израчунава вероватноћу да пример припада класи 1, на основу вредности улазних атрибута и научених коефицијената. Као и код линеарне регресије, неопходно је да нови подаци прођу исте трансформације као и подаци коришћени за тренирање.

Добијена вероватноћа се користи за доношење одлуке о класи. Најчешће се користи унапред дефинисан праг (на пример 0.5), али се у пракси овај праг може прилагођавати у зависности од пословних захтева и односа између алфа и бета грешке. На пример, у проблемима детекције преваре или медицинске дијагностике често се користи нижи праг како би се смањио број пропуштених случајева.

Поред саме класификације, вероватноће које модел даје могу се користити за рангирање или приоритизацију случајева. На пример, клијенти са већом вероватноћом одласка могу бити приоритетно обухваћени акцијама задржавања.

\subsection{Закључак}
Откривене законитости у форми асоцијативних правила, кластера, стабала одлучивања, регресионих модела или вештачких неуронских мрежа омогућавају организацијама да аутоматизују пословне процесе, скрате време извршења задатака и смање трошкове. При томе се мора имати на уму да није свако откривено правило законитост, нека правила могу пре да буду случајност, па се сви добијени модели обавезно евалуирају са статистичким тестовима.

Иако је ОЗП моћан алат, не представља замену за класичну статистичку анализу нити за стручно знање доносиоца одлука, већ њихову допуну. Процес ОЗП захтева активно учешће људи кроз све фазе — од разумевања пословног проблема и података, преко припреме и моделовања, до евалуације и примене. Будућност лежи у комбинацији две генерације система — оних који су вођени експертским знањем и оних који су засновани на подацима — где се интегришу велике количине доступних података и не занемарује експертско знање које се може уградити у моделе за поспешивање процеса доношења одлука.

\section{Литература}
\begin{itemize}
  \item Чупић М, Сукновић М (2008). \textit{Одлучивање}. Факултет организационих наука, Београд.
  \item Чупић М, Новаковић Т, Свилар М (1992). \textit{Генератори и апликације система за подршку одлучивању 1}. Научна књига, Београд.
  \item Делибашић Б, Сукновић М, Јовановић М (2010). \textit{Алгоритми машинског учења за откривање законитости у подацима}. Факултет организационих наука, Београд.
  \item Ester M, Kriegel H-P, Sander J, Xu X (1996). A Density-Based Algorithm for Discovering Clusters in Large Spatial Databases with Noise. 2nd International Conference on Knowledge Discovery and Data Mining.
  \item Fayyad UM, Irani KB (1992). The attribute selection problem in decision tree generation. Proceedings of the Tenth National Conference on Artificial Intelligence.
  \item Larose D (2004). \textit{Discovering Knowledge in Data, an Introduction to Data Mining}. John Wiley \& Sons.
\end{itemize}